\subsubsection{Entendendo vetores}

Ao abordar o conceito de vetores, entramos em um dos fundamentos da organização 
de dados em programação, especialmente na linguagem C. Vetores permitem agrupar 
dados do mesmo tipo em uma única estrutura acessível através de ponteiros. Isso 
simplifica a manipulação de coleções de dados como, por exemplo, um conjunto de 
idades tal qual no exemplo proposto. Vamos ver como isso pode ser implementado:

\begin{code}
    #include <stdio.h>
    
    int main(void)
    {
        // Declara um vetor de inteiros com espaço para 5 idades.
        int idades[5] = {23, 45, 56, 19, 30};
        int soma = 0;
        float media;

        // Calcula a soma das idades
        for (int i = 0; i < 5; i++) {
            soma += idades[i];
        }

        // Calcula a média das idades
        media = (float)soma / 5;

        printf("A média das idades é: %.2f\n", media);

        return 0;
    }
\end{code}

Neste exemplo, declaramos um vetor \texttt{idades} capaz de armazenar 5 valores 
inteiros, correspondentes às idades de pessoas aleatórias. Em seguida, 
utilizamos um \texttt{for loop} para somar todas as idades armazenadas no vetor. 
Após calcular a soma, dividimos esse valor pelo número total de idades para 
encontrar a média, que é armazenada na variável \texttt{media} e, por fim, 
exibida na tela.

Algumas nomenclaturas que podem facilitar o entendimento acerca dos vetores 
são as seguintes:

\begin{itemize}
        \renewcommand{\labelitemi}{\tiny$\blacksquare$}
    \item \textbf{Tamanho:} Refere-se à quantidade de espaço total alocada para o vetor 
        na memória, o tamanho de um vetor estático é definido durante a sua 
        declaração e \textbf{não pode ser alterado durante a execução do programa.}
        Geralmente o tamanho do vetor se encontra entre os colchetes durante a 
        declaração do mesmo ou é definido de maneira implícita pelo próprio compilador.
    \item \textbf{Comprimento ou quantidade de itens:} Refere-se à quantidade de elementos 
        atualmente armazenados no vetor, esse valor pode ser menor que o tamanho 
        total do vetor, mas nunca maior.
    \item \textbf{Índice:} Refere-se a posição de um item específico dentro do 
        vetor. Os índices são valores que podem ser utilizados como um número 
        de identificação para um item específico em um vetor de tamanho \(N\) 
        e podem ir de 0 até \(N\) - 1.
\end{itemize}

Um detalhe importante a se saber sobre o código acima, é que quando declaramos 
um vetor para guardar a idade de 5 pessoas, na verdade estamos alocando o 
espaço necessário para guardar 5 números, em sequência, na memória do computador. 
Ou seja, ao declarar:

\begin{code}
        int idades[5] = {23, 45, 56, 19, 30};
\end{code}

Estamos criando na memória do computador uma sequência de dados semelhante a seguinte:

\begin{figure}[htbp]
\centering
\includegraphics[scale=0.2]{array-in-memory-example}
\caption{Os valores do endereço da imagem acima são apenas ilustrativos.}
\label{fig:arrayInMemory}
\end{figure}

Ao considerarmos a criação de uma sequência de dados na memória do computador, 
conforme ilustrado anteriormente, compreendemos que um vetor é essencialmente 
uma série de valores organizados de maneira contínua na memória. A variável do 
vetor só precisa guardar o endereço do primeiro número na memória, semelhante 
ao que um ponteiro faria, e usa de um truque interessante para acessar cada um 
de seus elementos usando o valor do índice. No código acima, especificamente na 
seguinte linha:

\begin{code}
    soma += idades[i];
\end{code}

Usamos o valor de \(i\), que vai de 0 a 4 dentro do looping, como um valor de 
índice para indicar qual item queremos acessar dentro do vetor. Por exemplo, 
quando \(i\) é igual a 0, o computador acessa o endereço do primeiro número 
contido na variável do vetor e pula \(X\) itens a frente na memória baseado no 
valor passado entre colchetes.

\begin{postscript}
    Caso você queira entender como esse processo ocorre, de maneira mais 
    técnica, seria como se o compilador convertesse o trecho de código acima 
    para o trecho abaixo:

    \begin{code}
        soma += *(idades + sizeof(int) * i)
    \end{code}

    Neste contexto, \texttt{idades} atua como um ponteiro para o endereço do primeiro 
    elemento do vetor. A operação \texttt{sizeof(int) * i} calcula o deslocamento em 
    bytes necessários para alcançar o elemento do vetor relativo ao valor de \(i\), considerando o 
    tamanho de um inteiro (expresso por sizeof(int)). Somando este deslocamento ao 
    endereço base armazenado em idades, obtemos o endereço do elemento \(i\). 
    
    E em seguida usamos o operador de deferenciação * para acessar o valor 
    armazenado neste endereço específico, permitindo sua adição a soma.
\end{postscript}

\subsubsection{Utilizando vetores de texto}

O uso mais comum de vetores em C é para armazenar textos longos, como nomes e 
frases. Como a linguagem não possui tipos de variáveis nativos capazes de 
armazenar textos, nomes ou frases, recorre-se a vetores de variáveis do tipo 
\texttt{char} para representá-los. Vejamos o exemplo abaixo para ilustrar esse 
conceito:

\begin{code}
    #include <stdio.h>

    int main(void)
    {
        // Ao deixar o espaço entre cochetes vazio
        // o compilador calcula automaticamente o tamanho
        // necessário para guardar todos os itens do vetor
        char nome1[] = "Alan";
        char nome2[] = "Vinicius";
        char nome3[] = "Janiel";

        // A linguagem C prevê a impressão de
        // vetores de caracteres com o formatador '%s'
        printf("1º Nome: %s\n", nome1);
        printf("2º Nome: %s\n", nome1);
        printf("3º Nome: %s\n", nome1);

        return 0;
    }
\end{code}

No exemplo acima, criamos vetores de caracteres para armazenar os nomes de três 
pessoas, os quais são impressos posteriormente. Apesar de não existir um tipo 
de dado de texto nativo em C, a linguagem oferece implementações e convenções 
para a manipulação de vetores de caracteres como textos.

Um detalhe extremamente importante a mencionar é que, em relação ao uso de 
vetores para armazenar informações de texto, muitos projetos e bibliotecas 
utilizam um caractere especial para sinalizar o fim do vetor, conhecido como 
caractere nulo. Carácter esse que caso não esteja presente no vetor, ou até 
mesmo mal posicionado, pode ocasionar erros de acesso de memória inválida ou 
perda de informação.

Devido à capacidade de tamanho variável dos vetores de texto, não é possível 
determinar com certeza o tamanho do vetor quando ele é recebido por uma função, 
por exemplo. Para evitar erros relacionados ao acesso de memória não autorizada, 
uma convenção foi estabelecida: adicionar o carácter `\textbackslash0' ao final de cada 
vetor de texto, sinalizando o seu término.

\begin{code}
    char nome1[] = "Alan";
    char nome2[] = "Vinicius";
    char nome3[] = "Janiel";
\end{code}

No caso acima, quando declaramos e inicializamos os vetores de texto com os 
nomes, o compilador calcula o tamanho necessário para que o vetor possa conter 
todos os caracteres e mais um byte extra que deve conter o carácter nulo adicional.

\subsubsection{Declarando vetores de texto manualmente}

Outra forma não tão prática de definir um vetor, mas que ilustra bem a 
utilização do carácter nulo em um vetor, é a seguinte:

\begin{code}
    #include <stdio.h>

    int main(void)
    {
        char nome[5];
        nome[0] = 'A';
        nome[1] = 'l';
        nome[2] = 'a';
        nome[3] = 'n';
        nome[4] = '\0';

        printf("Nome: %s\n", nome);

        return 0;
    }
\end{code}

Aqui, declaramos manualmente um vetor com o tamanho de 5 caracteres e 
inicializamos cada um dos itens com os caracteres desejados seguidos ao 
carácter nulo especial no fim do vetor.
